\section{Introduction}
The number of exactly distinguishable posterior directions can jump when
an additive relation between observation exponents disappears. The
persistent memory needed at a fixed prediction accuracy need not jump:
new directions may have arbitrarily small amplitude. We study this
question for positive sparse monomial experiments, and determine how their sharp memory order is realized by finite
numerical programs. The numerical scale can itself be selected from
finite data, without first identifying an additive collision stratum.

Let $A=\{0=a_0<a_1<\cdots<a_{r-1}=D\}$, let
$W_A=\operatorname{span}\{t^a:a\in A\}$, and let $mA$ be the
set of sums of exactly $m$ members of $A$, with $0A=\{0\}$.
Write $h_A(m)=|mA|$. For the exact information statements the parameter
interval is $I=[l,u]$, $0\le l<u<\infty$, and the prior $\mu$ has
full support. For collision-uniform statements we take $I=[0,1]$.
A positive finite detector spans $W_A$, and a command in
$[\eta,1-\eta]^J$ retains or rejects each detector symbol. Every
reported symbol, including rejection, is counted as a trial. Its
likelihood is bounded below by a fixed positive number $\kappa$.
The precise experiment and common future-only equivalence are defined
in Section~\ref{sec:experiment}.

The one-step exponents range over a fixed compact subset $K$ of the
strictly ordered positive chamber. One fixed full-column-rank coefficient
matrix specifies the detector on $K$, with uniform positivity. Neither
semialgebraicity of $K$ nor a density for $\mu$ is assumed. The horizon
$N\ge2$ and prior are fixed. Choose a fixed
$H\ge\max\{1,N\max_Ka_{r-1}\}$. For each future length $m$, retain
all nonconstant formal homogeneous labels $\alpha\in\mathbb N_0^r$,
$|\alpha|=m$, including labels that have the same exponent. Their
normalized nodes are $x_\alpha=H^{-1}\alpha\cdot a$ and their number
is $q_m=\binom{m+r-1}{r-1}-1$. Set
\[
 \mathcal V_{m,\ell}(a)=\max_{|F|=\ell}
         \prod_{\{\alpha,\beta\}\subset F}|x_\alpha-x_\beta|,
 \qquad \mathcal V_{m,1}=1,
\]
where $F$ ranges over subsets of formal labels. Zero volumes remain
zero. Put $p_{n,m}=\min\{n(r-1),q_m\}$ and
\[
 \Psi_{n,m}(M,a)=\max_{1\le\ell\le p_{n,m}}
           (\mathcal V_{m,\ell}(a)/M)^{2/\ell},\qquad
 \Xi_N(M,a)=\max_{1\le n<N}\Psi_{n,N-n}(M,a).
\]
The checkpoint profile is zero at an endpoint. At each checkpoint the
query is selected independently after compression from a fixed finite
spanning menu of $m$-trial product probes. The excess squared loss is
the average squared error of its conditional event probabilities.
The notation $\overline R_{M;n,m}^a$ denotes optimal unconditional
checkpoint regret under independent uniform acquisition commands;
$R_{M;n,m}^{a,\max}$ denotes optimal worst-history regret.
Their causal counterparts $\mathfrak R^{a,\mathrm{av}}_{M,N}$ and
$\mathfrak R^{a,\max}_{M,N}$ minimize the maximum checkpoint criterion
over one stage-compatible filter. Only one checkpoint and its remaining
query trials are executed in a run.

\begin{theorem}[Sharp causal memory and certified numerical realization]
\label{thm:effective-main}
Under the preceding fixed-experiment hypotheses, with $r\ge2$, the
checkpoint and causal regret laws have orders $\Psi_{n,m}$ and $\Xi_N$,
respectively, uniformly in $a\in K$ and integers $M\ge1$.
Suppose the command cube is rationally specified and certified finite
approximations to the cell coefficients and prior moments through formal
degree $N$ are available at requested tolerances, together with coefficient and
positivity bounds.
An algorithm constructs integer transition tables and dyadic output
tables with at most $M$ persistent labels at each stage and maximum
worst-history regret at most $C\Xi_N(M,a)$. For every such filter with
a fixed memoryless command-coding rule, the maximum unconditional
checkpoint regret under the common exploration law is at least
$c\Xi_N(M,a)$.

The algorithm requires no collision stratum, additive-gap lower bound,
or value of $\Xi_N$ as input. Its selected numerical precision satisfies
\[
 b(M,a)=\tfrac12\log_2(1/\Xi_N(M,a))+O(1),
\]
and a sufficient read-only program size is
$O(M\Xi_N(M,a)^{-J/2}\log(M+1))$ bits. In particular, setting
$d_0=(r-1)\lfloor N/2\rfloor$, one may uniformly use
\[
 b(M)\le d_0^{-1}\log_2(M+1)+O(1),\qquad
 P(M)=O(M^{1+J/d_0}\log(M+1)).
\]
These numerical bounds are sufficient; no precision or program-length
converse is asserted. The running filter retains only its index, not a
posterior vector, historical tape, or numerical oracle. For computable
numerical inputs the synthesis is uniform and terminating. All constants
may depend on the fixed prior and horizon, but not on $a$ or $M$.
\end{theorem}

For $r=1$, all report factors are constant in the hidden parameter; the
posterior-dependent state is a singleton and one label suffices. The
exact information classification below includes this case.

The geometric assertions are proved in
Theorems~\ref{thm:intrinsic-checkpoint} and
\ref{thm:intrinsic-streaming}. The finite-table construction is given
in Section~\ref{sec:effective}; its stronger resource conclusions and
adaptive stopping rule are proved in
Theorem~\ref{thm:profile-adaptive} and
Corollary~\ref{cor:uniform-resources}. The earlier fine-mesh bounds
remain valid consequences and are recorded with their proofs. For an
arbitrary, possibly noncomputable full-support prior, the intrinsic
classification is unchanged. The effective statement uses supplied
numerical advice, rather than an algorithm for obtaining unsupplied
noncomputable moments. With an external clock the changing state has
$M$ values; storing the clock costs a factor depending only on $N$.
The program and temporary workspace are accounted for separately.

\subsection{Structure of the proof}
The first step is a pairing between the attainable product tangent and
the complete future test flag. At an interior binomial tuple, the product
tangent contains $n(r-1)+1$ distinct monomials. Strict mixed-moment
positivity holds for every full-support prior. Complete Hermite blocks
retain this positivity at exact collisions; normalization removes exactly
one rank. The local image carries an unconditional volume minorization
which includes the probability of acquiring the failure prefix. A
separate bounded-format argument covers the entire reachable image with
the attainable dimension truncated at $p_{n,m}$.

The determinant profile has a classical finite spectral meaning:
Proposition~\ref{prop:exterior-profile} identifies it, up to fixed
constants, with the exterior singular-value profile of a Vandermonde
matrix. That statement alone does not imply statistical attainment.
The arbitrary-prefix Hermite statement in
Lemma~\ref{lem:prefix-hermite} is likewise classical
\cite[Proposition 7]{deBoor}. The normalized product, acquired-history
probability, and compatible transitions provide the additional
statistical structure.

The numerical construction never uses singular coordinates. It forms
finite lists of feasible histories, evaluates their raw states to a
certified accuracy, and applies classical farthest-first selection
\cite{Gonzalez}. Every transition is evaluated by appending a report and
command to an actual representative history. Approximate vectors outside
the reachable set are not passed to a posterior update. Positive evidence
certifies evaluation from a finite moment table even when coincident
formal labels receive different approximate values.

The final step makes the numerical scale observable. If $r$ is the
computed covering radius of an approximate history list with mesh and
evaluation tolerance $h$, and $e$ is the largest true $M$-centre covering
radius, then
\[
 \max\{0,r/2-h\}\le e\le r+(A+2)h.
\]
The rule $r\ge4(A+2)h$ therefore stops at $h\asymp e$ without knowing
$e$. Attainable geometry identifies $e^2\asymp\Xi_N$. This gives a
finite-data interpretation of the full profile, not only of one of its
power-law floors. A separated future chain gives the uniform precision
bound. Transition and decoder residuals provide separate finite checks
on the program actually returned.

For the intersecting two-parameter family, the resulting program-size
phases track the full three-term regret profile. High-order contact may
place a nonzero exponent gap far below the chosen mesh, without changing
the guarantee. The algorithm does not identify the contact order or
collision tree. The stable common-advice theorem is then proved in
Section~\ref{sec:construction-stability}. Complete observation
algebra, affine, five- and seven-trial, decision, and mechanical proofs
are given in the appendices. Their separate assumptions are retained.


\subsection{Imperfect numerical descriptions}
The geometric law has a further consequence when a finite moment name
is compatible with more than one experiment. On a dominated family of
full-support priors, the constants remain uniform. Pairing identical
histories proves Hausdorff stability of the attainable state sets in
the raw coordinates; optimal covering radii are therefore Lipschitz
in the finite moment data. A common-name stopping rule yields one
$M$-label program, not a family selected using an unsupplied true model,
with regret $O(\Xi_N(M,a)+\delta^2)$ in every compatible experiment.
These assertions and their precise uniformity hypotheses are stated
in Lemma~\ref{lem:prior-envelope},
Theorem~\ref{thm:common-advice}, and
Corollary~\ref{cor:uncertain-monomial}.

The extra term has a matching elementary lower witness. A fixed positive
rank-two detector admits two full-support priors with the same
$\delta$-accurate finite moment name and with a physical prediction gap
of order $\delta$ after an identical positive-probability acquisition
event. Proposition~\ref{prop:advice-floor} gives regret of order
$M^{-2}+\delta^2$ for this experiment. The two-point lower-bound method
is elementary; the assertion being added is its compatibility with the
same physical, budgeted prediction problem and the collision-uniform
upper construction. It is not a claim that every finite-precision
implementation has this uncertainty model.
